\documentclass[12pt]{article}
\usepackage{amsmath}
\usepackage{hhline}

\begin{document}

Die kubische Gleichung $x^3 + rx^2 + sx + t = 0$ kann mittels der Substitution
$y = x + \frac{r}{3}$ in die sogenannte reduzierte Gleichung $y^3 + py + q = 0$
überführt werden. Die Wurzeln $y_1$, $y_2$ und $y_3$ der reduzierten Gleichung
lassen sich gemäß folgender Tabelle mit
\[
  R := (\operatorname{sgn} q)\sqrt{\frac{|p|}{3}}
  \quad \text{und} \quad
  D := \left(\frac{p}{3}\right)^{\!3} + \left(\frac{q}{2}\right)^{\!2}
\]
bestimmen:

\renewcommand{\arraystretch}{1.6}
\begin{center}
\begin{tabular}{c || c | c | c}
  \multicolumn{1}{c}{} & \multicolumn{2}{c|}{$p < 0$} & \\
  \cline{2-3}
  \multicolumn{1}{c}{} & $D \leq 0$ & $D > 0$ & $p > 0$ \\
  \hhline{~|=|=|=}
  \multicolumn{1}{c|}{} & $\cos\varphi = \dfrac{q}{2R^3}$
  & $\cosh\varphi = \dfrac{q}{2R^3}$
  & $\sinh\varphi = \dfrac{q}{2R^3}$ \\
  \hhline{~|=|=|=}
  $y_1$ & $-2R\cos\dfrac{\varphi}{3}$
        & $-2R\cosh\dfrac{\varphi}{3}$
        & $-2R\sinh\dfrac{\varphi}{3}$ \\
  \hline
  $y_2$ & $-2R\cos\!\left(\dfrac{\varphi}{3}+\dfrac{2\pi}{3}\right)$
        & $R\cosh\dfrac{\varphi}{3}+i\sqrt{3}\,R\sinh\dfrac{\varphi}{3}$
        & $R\sinh\dfrac{\varphi}{3}+i\sqrt{3}\,R\cosh\dfrac{\varphi}{3}$ \\
  \hline
  $y_3$ & $-2R\cos\!\left(\dfrac{\varphi}{3}+\dfrac{4\pi}{3}\right)$
        & $R\cosh\dfrac{\varphi}{3}-i\sqrt{3}\,R\sinh\dfrac{\varphi}{3}$
        & $R\sinh\dfrac{\varphi}{3}-i\sqrt{3}\,R\cosh\dfrac{\varphi}{3}$ \\
\end{tabular}
\end{center}

\end{document}